\documentclass[minitoc, hidelinks, headlogo, framedcover]{tfeEPCC}
\usepackage[spanish, es-tabla]{babel} % Comentar si se va a realizar en inglés
\usepackage{booktabs}   % reglas horizontales pro
\usepackage{float}      % permite usar [H] en tablas
%\usepackage{parskip} % Quitar comentario si se quieren distinguir los párrafos con saltos y quitar las sangrías

\usepackage{lipsum} % Este paquete es para poner texto genérico en los capítulos, quitar cuando ya se hayan redactado
\usepackage{tcolorbox}
\usepackage{xcolor}
\usepackage{listings}
\usepackage{longtable}
\usepackage{array}

% Rellenar con los datos del estudiante y el trabajo
\estudiante{Ignacio Alcalde Torrescusa}{Grado}{INGENIERÍA INFORMÁTICA DEL SOFTWARE}
\title{ANÁLISIS DE LAS COMUNICACIONES Y CIBERSEGURIDAD EN LA ERA POST-CUÁNTICA}
\keywords{Palabras más importantes o descriptivas\and separadas por comas\and para facilitar la búsqueda si se deposita en repositorio}

\begin{document}
\sloppy 
\pagebreak 

%%%%%%%%%%%%% PORTADA %%%%%%%%%%%%%
% Poner convocatoria y año de presentación
\portadaTFE{Septiembre, 2025}

\pagestyle{empty}


%%%%%%%%%% CONTRAPORTADA %%%%%%%%%%
% Campos, separar con comas si se ponen varios: {
% studentGender
% tutor (obligatorio)
% tutorGender
% cotutor
% cotutorGender }
% Especificar un co-tutor lo añade a la contraportada
% Poner la letra f en cualquier campo de género añade una 'a' al título correspondiente (por ejemplo, studentGender = f hace que se muestre Autora)
\contraportadaTFE{tutor = Jesús Manuel Calle Cancho}

%%%% CAPÍTULOS INTRODUCTORIOS %%%%
\pagenumbering{gobble}
%\chapter*{Agradecimientos} % Apartado opcional, descomentar si se quiere añadir
% Dedicatorias del trabajo, pueden ser en español o inglés dependiendo de lo que se considere más apropiado

\chapter*{Resumen}
Este Trabajo de Fin de Grado presenta una revisión sistemática de la literatura (SLR) centrada en dos campos tecnológicos emergentes estrechamente relacionados como son: las comunicaciones cuánticas y la ciberseguridad en la era post-cuántica. El estudio examina de manera conjunta el estado del arte de cada disciplina, identificando avances y retos, así como el papel que desempeñarán en la seguridad de la información,  transformación de las redes y telecomunicaciones durante las próximas décadas.

En este contexto, surgen dos líneas de investigación complementarias. Por un lado, las comunicaciones cuánticas, orientadas a la creación de canales de transmisión seguros bajo principios físicos como la distribución cuántica de claves (QKD) y por otro, la ciberseguridad post-cuántica, enfocada en el desarrollo de algoritmos criptográficos resistentes ante ataques de computadoras cuánticas.

La revisión sistemática se ha realizado de manera rigurosa y totalmente transparente, se han empleado bases de datos científicas de alto impacto y herramientas especializadas para garantizar la trazabilidad y calidad del análisis. Este enfoque ha permitido ofrecer una visión estructurada y actualizada de cada área, detectar tendencias, tecnologías clave e identificar vacíos en la literatura, además de proponer nuevas líneas de investigación futuras.

Los resultados ponen de manifiesto la urgencia de iniciar la transición hacia infraestructuras criptográficas resistentes a la computación cuántica y el desarrollo de estándares internacionales que garanticen comunicaciones y sistemas seguros en un entorno tecnológico que a día de hoy, ya se encuentra bajo una amenaza real.

% Índices con enlaces de las secciones, tablas y figuras
\newpage
\pagestyle{plain}
\pagenumbering{Roman}
\setcounter{page}{1}
\tableofcontents
%\pagebreak
\newpage
\listoftables
\newpage
\listoffigures
\newpage


%%%%%%% ESTRUCTURA DEL TFG %%%%%%%%%%%%%%
% A.PORTADA (según la estructura indicada a continuación) 
% B.CONTRAPORTADA (según la estructura indicada a continuación) 
% C.ÍNDICE GENERAL DE CONTENIDOS  
% D.ÍNDICE DE TABLAS  
% E.ÍNDICE DE FIGURAS 
% F.RESUMEN (Podrá incluirse también en inglés, si así lo indica el Tutor1) 
% G.CUERPO DEL TRABAJO (según la estr
% uctura indicada a continuación) 
% H.REFERENCIAS BIBLIOGRÁFICAS 
% (Según norma ISO690)
% I.ANEXOS, si los hubiera
%%%%%%%%%%%%%%%%%%%%%%%%%%%%%%%%%%%%%%%
%%%%%% Cuerpo del trabajo
% 1.INTRODUCCIÓN 
% 2.OBJETIVOS 
% 3.ANTECEDENTES / ESTADO DEL ARTE 
% 4.MÉTODOLOGÍA 
% 5.IMPLEMENTACIÓN Y DESARROLLO (Cuando proceda) 
% 6.RESULTADOS Y DISCUSIÓN 
% 7.CONCLUSIONES 
%%%%%%%%%%%%%%%%%%%%%%%%%%%%%%%%%%%%%%%%%%%
%%%%%%%%%%%%%%%%%%%%%%%%%%%%%%%%%%%%%%%%%%%
% Tamaño: Normalizado UNE A-4, salvo planos. 
% Tipo y tamaño de letra del texto: Times New Roman 12 pt, Arial 12 pt o similar. 
% Interlineado del texto: 1,5 líneas. 
% Márgenes del texto: Superior, Inferior y Derecha, 2,5 cm; Izquierda, 4 cm. 
% Numeración de páginas en margen inferior derecha y tamaño 8 pt. 
% Las  figuras  serán  numeradas  y  tituladas  debajo  de  las  mismas  (indicando  su  fuente  si  no  son  de  elaboración  
% propia). 
% Las  tablas  serán  numeradas  y  tituladas  encima  de  las  mismas  (indicando  su  fuente  si  no  son  de  elaboración  
% propia). 
%%%%%%%%%%%%%%%%%%%%%%%%%%%%%%%%%%%%%%%%%%%
%%%%%%%%%%%%%%%%%%%%%%%%%%%%%%%%%%%%%%%%%%%
%%%%%%%%%%%%%%%%%%%%%%%%%%%%%%%%%%%%%%%%%%%
%%%%%%%%%%%%%%%%%%%%%%%%%%%%%%%%%%%%%%%%%%%


\pagenumbering{arabic}
\pagestyle{fancy}
\setcounter{page}{1}

%%%%%%%% INTRODUCCIÓN %%%%%%%%
\import{chapters}{01_introduccion.tex}

%%%%%%%% OBJETIVOS %%%%%%%%
\import{chapters}{02_objetivos.tex}

%%%%%%%% ESTADO DEL ARTE %%%%%%%%
\import{chapters}{04_estado_del_arte.tex}

%%%%%%%% IMPLEMENTACIÓN Y DESARROLLO %%%%%%%%
\import{chapters}{05_materiales.tex}

%%%%%%%% SLR CIBERSEGURIDAD EN LA ERA POST-CUÁNTICA %%%%%%%%
\import{chapters}{06_SLR_CiberseguridadPostCuántica.tex}

%%%%%%%% SLR CIBERSEGURIDAD EN LA ERA POST-CUÁNTICA %%%%%%%%
\import{chapters}{07_SLR_ComunicacionesCuánticas.tex}

%%%%%%%% CONCLUSIONES %%%%%%%%
\import{chapters}{08_Conclusiones.tex}


% Si no se van a incorporar anexos, comentar el comando '\startappendices' junto a todos los imports de la carpeta 'appendices'
%%%%%%%%%%%%%%%%%%%%%%%%%%%%%%%%%%
%%%%%%%%%%% ANEXOS %%%%%%%%%%%%%%%
%%%%%%%%%%%%%%%%%%%%%%%%%%%%%%%%%%
\startappendices

%%%%%%%% ANEXO 1 %%%%%%%%
\import{appendices}{01_resultados_slr.tex}

%%%%%%%% ANEXO 2 %%%%%%%%
\import{appendices}{02_cuadernos_colab.tex}

\pagebreak
%%%%%%%%%%%%%%%%%%%%%%%%%%%%%%%%%%
%%%%%%%%%% AL FINAL %%%%%%%%%%%%%%
%%%%%%%%%%%%%%%%%%%%%%%%%%%%%%%%%%
\pagestyle{empty}
%%%% https://en.wikibooks.org/wiki/LaTeX/Bibliography_Management
%%%% https://www.bibtex.com
%%%%%%%% BIBLIOGRAFÍA %%%%%%%%
% Se puede cambiar el estilo de la bibliografía, por defecto es 'unsrturl'
% Por ejemplo: \makebibliography{alphaurl}
\makebibliography
\end{document}